\section{Conclusion}\label{sec:conclusion}

The work produced five domain-general principles and the 106 numbered clauses of the Specification, which turn them into obligations a project can be checked against \citep{tbf2026,pss2026}. The clauses fix 24 required records and their relations. Each project declares its storage in a complete local path map, and optional software and hardware profiles supply filesystem examples.

In this work, mechanisms already used in practice became rules written to apply across domains, through a conversion in which every step keeps the record of the step before it. That answers the first of the two questions stated in the introduction. The conversion ran from 219 recorded practice statements through 79 domain-general forms to the clauses, each stage under its own control and with a recorded decision for every candidate.

What keeps the chain inspectable rather than asserted is the bracket at the end of every clause. Its links close in both directions: no clause stands without a stated basis, and no adopted form was collected and left unused.

Recorded grounds and relations answer the second question by separating three things a flat document holds together: missing support, conflicting requirements, and the reach of a proposed change.

Recorded grounding is complete and purpose is not. Most of the 25 clauses without a recorded path to a stated goal govern how the project record is organised rather than what an engineered system must do.

The conflicts are a property of the representation. All 39 flagged pairs trace to one missing slot in the rule schema, which cannot yet hold the restricting part of a prescription.

The reach of a change comes back as a partition. A proposed edit to the clause that mandates the minimum check set names 185 source nodes at the last depth reached only by an asserted relation or an explicit reference. It also proves 3{,}090 of the 3{,}853 source nodes unreachable from that clause at any depth.

The relation graph rebuilds from its recorded decisions to the same logical fingerprint, and the rebuilt file is byte-identical to the recorded one. Repeatability from recorded decisions is what that establishes, over four documents whose corpus-scale relations are delegated-AI single-lane output rather than human ratification. Whether the purpose is the right purpose, and whether the residual risk is tolerable, stay with the responsible human.

The two requirements stated in the introduction are carried by the Specification on itself, as a grounding layer complete in both directions and a traversal that names what an edit reaches. One relation falls short. Those 25 clauses are requirements whose purpose the document does not record, which is the thing it asks every project to record.

Four parts of the framework are named in its goal document and not built: the AI-native execution methods, the domain playbooks, the onboarding path, and the templates and checklists \citep{goal2026}. None of them is reported in this paper.

Independent application is the next measurement of the Specification. Until it runs, the evidence stated in this paper is internal to the work: the derivation is recorded, the methods and findings are reported with their review conditions, and its effect on a project is unmeasured.
